% =========================================================
% Proof Stub: epsilon-Characterization of Supremum and Infimum
% Source: volume-iii/analysis/bounding/notes/notes-bounds-extremal-values.tex
% Mode: standard
% =========================================================

\clearpage
\phantomsection
\hypertarget{proof-RA-ABB-P-eps-char-sup}{}

\subsubsection[$\varepsilon$-Characterization of Supremum]{Proof --- $\varepsilon$-Characterization of Supremum and Infimum}

\bigskip

\noindent
\textbf{Source.}
Abbott, \emph{Understanding Analysis}, Chapter~1; see also notes
\texttt{notes-bounds-extremal-values.tex}, Proposition~(\ref{prop:eps-char}).

\vspace{0.75em}

\noindent
\textbf{Goal.}
Show that $s = \sup A$ if and only if (a) $s$ is an upper bound for $A$,
and (b) for every $\varepsilon > 0$ there exists $a \in A$ with
$s - \varepsilon < a$.

\vspace{0.75em}

\noindent
\textbf{Logical form.}
\[
  s = \sup A
  \;\iff\;
  \Bigl[\,
    (\forall a \in A,\; a \le s)
    \;\wedge\;
    (\forall \varepsilon > 0,\; \exists a \in A,\; s - \varepsilon < a)
  \,\Bigr].
\]

\vspace{0.75em}

\noindent
\textbf{Proof strategy.}
This is a biconditional. For ($\Rightarrow$): the upper-bound clause is
immediate; for the $\varepsilon$-clause, if no such $a$ existed then
$s - \varepsilon$ would be a smaller upper bound, contradicting $s$ being
least. For ($\Leftarrow$): if $u$ is any upper bound, suppose for
contradiction $s > u$; then $\varepsilon := s - u > 0$ and no $a > s - \varepsilon = u$ exists in $A$, contradicting (b).

\vspace{0.75em}

\noindent
\textbf{Proof.}
\begin{proof}
\Claim $s = \sup A \iff$ (a) $s$ is an upper bound for $A$, and (b)
$\forall \varepsilon > 0,\, \exists a \in A$ with $s - \varepsilon < a$.

\Given $A \subseteq \mathbb{R}$ nonempty and bounded above, $s \in \mathbb{R}$.

\Goal Prove the biconditional.

\Strategy ($\Rightarrow$) Use the definition of $\sup$ and argue by
contradiction for the $\varepsilon$-clause.
($\Leftarrow$) Use (b) to show no smaller upper bound can exist.

\medskip

% --- ($\Rightarrow$) ---

% --- ($\Leftarrow$) ---

\end{proof}

\begin{remark*}[Proof shape]
The $(\Rightarrow)$ direction uses proof by contradiction for the
$\varepsilon$-clause. The $(\Leftarrow)$ direction is a direct argument
that $s$ beats every upper bound.
\end{remark*}

\begin{remark*}[Dependencies]
\begin{itemize}
  \item Definition of supremum (least upper bound).
  \item Trichotomy of $\mathbb{R}$: for any $s, u \in \mathbb{R}$, exactly
        one of $s < u$, $s = u$, $s > u$ holds.
\end{itemize}
\end{remark*}
